% Study sheet: Volume I, Chapter 2, System Paradigms.
% Numbers: \val{...} from system_paradigms_studysheet_numbers.tex, written by system_paradigms_studysheet_numbers.py.
\documentclass{studysheet}
\input{system_paradigms_studysheet_numbers}
\sheetvolume{1}
\sheetchapter{2}{System Paradigms}
\sheetsources{Numbers from the mlsysim registry and the chapter's scenarios; questions marked assumed use new inputs.}

\begin{document}
% ===================================================================== FRONT: LEARN
\frontside
\begin{bigidea}
Where a model runs decides which \hl{physical limit} it meets first: light, heat or memory. Cloud,
edge, mobile and TinyML are regimes set by those limits, from megawatts to milliwatts. The task is to
\hl{find the limiting constraint}, not the fastest machine.
\end{bigidea}

\begin{leftcol}
\sheetsec{Key ideas}
\begin{ideas}
\idea{Three walls.}{Light sets a round-trip floor; heat caps sustained power; memory makes moving
  data a growing share of the cost.}
\idea{Silicon contract.}{Bandwidth, compute rate and overhead are the iron law's terms; memory
  capacity and energy are preconditions, checked first.}
\idea{Four archetypes.}{Two name an iron-law term, two a precondition. Batching reuses each loaded
  weight, moving a model from bandwidth- toward compute-bound.}
\idea{Bounds reject.}{A remote path whose lower bound misses the deadline is out; passing only makes
  it a candidate to measure.}
\idea{Screen, then cost.}{First: does it need ML? Then privacy, latency, resources, then lifecycle
  cost among survivors. Excluded tiers stay out.}
\idea{Else, go hybrid.}{Train-serve split; hierarchical (a local filter wakes a bigger tier);
  progressive (a variant per tier). Keep preprocessing and versions in sync.}
\end{ideas}

\vspace{10pt}
\sheetsec{What limits it first?}
{\centering
\begin{tikzpicture}[sheet, every node/.style={align=center}]
  \tikzset{q/.style={nbox, text width=104pt, minimum height=20pt},
           o/.style={accbox, text width=96pt, minimum height=20pt}}
  \node[q] (q1) at (60,90) {Does the working set fit in memory?};
  \node[q] (q2) at (60,55) {Can the energy budget pay for the run rate?};
  \node[q] (q3) at (60,20) {Is the compute term larger than the data term?};
  \node[o] (o1) at (190,90) {\textbf{capacity-bound}\\ DLRM $\cdot$ cloud};
  \node[o] (o2) at (190,55) {\textbf{power-bound}\\ KWS $\cdot$ TinyML};
  \node[o] (o3) at (190,20) {\textbf{compute-bound}\\ ResNet-50 training $\cdot$ cloud};
  \node[o] (o4) at (190,-13) {\textbf{bandwidth-bound}\\ GPT-2 decode $\cdot$ cloud};
  \draw[arr] (q1) -- node[small, above] {no} (o1);
  \draw[arr] (q2) -- node[small, above] {no} (o2);
  \draw[arr] (q3) -- node[small, above] {yes} (o3);
  \draw[arr] (q1) -- node[small, right, xshift=2pt] {yes} (q2);
  \draw[arr] (q2) -- node[small, right, xshift=2pt] {yes} (q3);
  \draw[arr] (q3.south) |- node[small, pos=0.25, right, xshift=2pt] {no} (o4.west);
\end{tikzpicture}\par}

\vspace{10pt}
\sheetsec{Four tiers}
\begin{sheettable}{@{}l c c L@{}}
\thead{Tier} & \thead{Power} & \thead{Memory} & \thead{Binds first}\\ \midrule
Cloud & MW & TB & compute or bandwidth; distance\\
Edge & $\sim$100 W & GB & fixed site capacity\\
Mobile & \val{tdp} W sustained & GB & energy: battery and heat\\
TinyML & mW & KB (SRAM) & memory capacity, then energy\\
\end{sheettable}
\end{leftcol}%
\begin{rightcol}
\sheetsec{Formulas you will use}
\fcard{Iron law: one request, and a stream}%
  {\begin{gathered}\Time = \frac{\DataVolume}{\Bandwidth} + \frac{\Operations}{\PeakRate\,\EfficiencyHw} + \Latency\\[2pt]
   \PipelinedExecutionTime \ge \max\!\Big(\frac{\DataVolume}{\Bandwidth},\ \frac{\Operations}{\PeakRate\,\EfficiencyHw},\ \IoTransferTime\Big) + \Latency\end{gathered}}%
  {One request pays the sum ($\PeakRate\,\EfficiencyHw$: achieved rate). A fully overlapped stream pays the
   slowest stage per task ($\IoTransferTime$: network or storage transfer).}%
  {Speeding a stage that is not the slowest helps one request, not the stream.}{}
\fcard{Light barrier}%
  {\MinimumRoundTripLatency = \frac{2 \times \text{Distance}}{\FiberPropagationSpeed},\qquad
   \FiberPropagationSpeed \approx \val{cfiber}\ \text{km/s}}%
  {About \val{per100} ms per 100 km, paid before any computing; faster servers cannot shorten it.}{}{}
\fcard{Data locality invariant}%
  {\RemotePathTime \ge \frac{\DataVolume}{\NetworkBandwidth} + \NetworkLatency
   + \frac{\Operations}{\RemotePeakRate\,\RemoteHardwareEfficiency}}%
  {Send, cross the network, compute. Above the deadline $\CompletionDeadline$: reject. Below: a candidate.}%
  {The largest input that passes may use only the time left after $\NetworkLatency$ and remote compute.}{}
\fcard{Energy on a battery}%
  {\CloudOffloadEnergy \approx \DataVolume \times \TransmissionEnergyPerByte \quad\text{vs.}\quad \LocalInferenceEnergy}%
  {Radio energy per byte sent (Wi-Fi or LTE: about \val{radio} mJ per MB) against one local
   inference; a faster cloud refunds none.}%
  {Battery caps total energy; heat caps sustained power (\val{tdp} W on a phone). At its temperature
   trip point the chip throttles, at any charge.}{}
\end{rightcol}

% ===================================================================== BACK: PRACTICE
\backside
\begin{answers}
\ans{Finish}{\val{ytx} $+$ \val{drresp} $=$ \val{ytot} ms $<$ \val{ydl} ms: a candidate to measure.
  \val{ydl} $-$ \val{drresp} $=$ \val{yleft} ms of sending carries \val{ycross} MB. The drone frame still
  takes \val{drtx} s to send: no server saves it (locality card).}
\ans{1}{Data \val{qtd} ms $>$ compute \val{qtc} ms: bandwidth-bound. Request \val{qser} ms, stream
  \val{qpipe} ms; faster math: \val{qser2} ms and still \val{qpipe2} ms (iron law card).}
\ans{2}{(a) TinyML, power-bound. (b) Edge: privacy, fixed capacity. (c) Cloud, compute-bound.
  Hierarchical plus train-serve split; sync preprocessing and versions (hybrid idea).}
\columnbreak
\ans{3}{Privacy removes the cloud (its \val{frtt} ms floor is moot). Resources remove TinyML (KB, not
  \val{fgb} GB) and mobile (\val{fw} W vs.\ \val{tdp} W). Edge: measure, then compare designs on
  lifecycle cost (tier table).}
\ans{4}{(a) \val{wmb} $\times$ \val{radio} $=$ \val{wtx} mJ, \val{wx}$\times$ local; local wins above
  \val{wbe} KB, however fast the cloud. (b) Heat: clock speed falls as chip temperature hits the trip
  point, at full charge (energy card).}
\end{answers}

\begin{leftcol}
\sheetsec{Worked example}
\begin{practicebox}\boxtitle{Can a drone offload obstacle detection to the cloud?}
Each \val{drmb} MB 4K frame goes over a \val{drmbit} Mb/s uplink; the network adds \val{drnet} ms, the
server \val{drcmp} ms. A result is due before the next frame: \val{drdl} ms at \val{drfps} frames/s.\par\vspace{6pt}
{\centering
\begin{tikzpicture}[sheet]
  \pgfmathsetmacro\sc{196/(\val{drtx}*1000+\val{drresp})}
  \pgfmathsetmacro\xa{\val{drtx}*1000*\sc}
  \pgfmathsetmacro\xb{\xa+\val{drresp}*\sc}
  \pgfmathsetmacro\xd{\val{drdl}*\sc}
  \draw[hotbox] (0,0) rectangle (\xa,16);
  \node[lab, text=white] at ({\xa/2},8) {send the frame \val{drtx} s};
  \draw[accbox] (\xa,0) rectangle (\xb,16);
  \node[lab, anchor=west] at ({\xb+3},8) {\val{drtot} s};
  \node[note, anchor=north east] at (\xb,-3) {network $+$ server \val{drresp} ms};
  \draw[ink, dashed, line width=0.7pt] (\xd,16) -- (\xd,24);
  \draw[white, line width=0.9pt] (\xd,0.6) -- (\xd,15.4);
  \node[note, anchor=south west] at (\xd-1,25) {deadline \val{drdl} ms};
\end{tikzpicture}\par}\vspace{4pt}
\begin{steps}
\item Send: \val{drmb} MB $\times$ \val{bits} bits/byte $\div$ \val{drmbit} Mb/s $=$ \res{\val{drtx} s}.
\item Lower bound: \val{drtx} s $+$ \val{drresp} ms $=$ \res{\val{drtot} s}, \val{drx}$\times$ the \val{drdl} ms deadline.
\item Even an empty input misses: \val{drresp} ms $>$ \val{drdl} ms. Reject; check the local path.
\end{steps}
\takeaway{A bound above the deadline rejects, whatever the server.}
\end{practicebox}

\vspace{12pt}
\sheetsec{Your turn: finish it}
\begin{practicebox}\boxtitle{Is the cloud a candidate for photo tagging?}
\assumed{} the same uplink and \val{drresp} ms; a \val{ymb} MB photo due within \val{ydl} ms.\par\vspace{4pt}
\textbf{Given:} send $=$ \val{ymb} MB $\times$ \val{bits} bits/byte $\div$ \val{drmbit} Mb/s $=$ \val{ytx} ms.\par
\textbf{You:} find the bound and what passing means; the largest photo that passes; whether a
faster server would save the drone frame.
\end{practicebox}
\end{leftcol}%
\begin{rightcol}
\sheetsec{Common mistakes}
\mistake{``The cloud has the most compute, so it answers fastest.''}%
  {A robot arm's stop served \val{armkm} km away: 2 $\times$ \val{armkm} km $\div$ \val{cfiber} km/s $=$
   \val{armrtt} ms of round trip, over its \val{armbudget} ms budget.}{}
\mistake{``A redesign drawing \val{thk}$\times$ less power leaves plenty of thermal headroom.''}%
  {A \val{thermw} W model throttles within \val{thtrip} s. Cutting power \val{thk}$\times$ reaches \val{thnew} W,
   exactly the \val{thceil} W ceiling: no headroom left, and no optimization raises the ceiling.}{}
\mistake{``Edge hardware is cheap, so edge is cheaper.''}%
  {\$\val{tcohw}/month of hardware vs.\ \$\val{tcocloud} of cloud compute, but network, maintenance and reliability work
   (\$\val{tconet} $+$ \$\val{tcomaint} $+$ \$\val{tcorel}) bring edge to \$\val{tcoedge}, \val{tcox}$\times$ the cloud.}{}

\vspace{2pt}
\sheetsec{Check yourself}
\begin{qs}
\item \assumed{} \val{qmb} MB and \val{qgf} GFLOP per inference; \val{qbw} GB/s, \val{qtf} TFLOP/s achieved,
  \val{qL} ms overhead. Which term limits it? Time per request, and per task in a fully overlapped
  stream (no I/O)? What does \val{qk}$\times$ faster math change?
\item A home camera: (a) a coin-cell motion detector lasting years; (b) a hub that identifies people,
  video kept at home; (c) monthly retraining. Each stage's tier and limit? Which hybrid patterns? What
  stays in sync?
\item \assumed{} inspection within \val{fdl} ms; raw images stay on site; the detector needs \val{fgb} GB
  and \val{fw} W; the cloud is \val{fkm} km away. Run the four checks. Which tiers survive? What is left?
\item \assumed{} a wearable radios a \val{wmb} MB clip, or infers locally for \val{wloc} mJ. (a) Which
  uses less energy, even with an infinitely fast cloud? (b) Its phone app slows after a minute at full
  charge: battery or heat? What would you watch?
\end{qs}
\end{rightcol}
\end{document}
